\chapter*{Abstract}
\addcontentsline{toc}{chapter}{Abstract}

Managing large-scale Capital Expenditure (CAPEX) portfolios requires the coordinated execution of numerous parallel projects with complex interdependency structures. A fundamental challenge in this context is the lack of transparency concerning specific skill requirements and their temporal distribution across the project lifecycle. Without precise, forward-looking insights into which competencies are required at which point in time, organizations are unable to proactively ensure on-time project delivery. Existing manual planning methods frequently fail to identify critical resource bottlenecks sufficiently early, resulting in project delays, cost overruns, and suboptimal investment outcomes.

This thesis develops a generalizable, application-oriented framework for an AI-supported recommendation engine designed to enhance strategic resource management in CAPEX portfolio environments. The framework addresses three interrelated challenges: (1) \textit{temporal skill-demand forecasting}, which utilizes advanced analytical models to predict the specific skill sets and capacities required across different project phases; (2) \textit{proactive bottleneck identification}, which detects potential resource scarcities early enough to enable timely strategic adjustments and risk mitigation; and (3) \textit{optimized allocation recommendations}, which suggest the most efficient distribution of available resources to meet project timelines while respecting organizational constraints.

The proposed system integrates machine learning techniques for demand forecasting with optimization algorithms for allocation planning, embedded in a decision-support architecture. Evaluation results demonstrate that \ldots % TODO: fill in after evaluation

\bigskip
\noindent\textbf{Keywords:} CAPEX portfolio management, resource allocation, skill-demand forecasting, bottleneck detection, AI recommendation engine, optimization
